\documentclass[11pt]{article}
\usepackage{amsmath,amssymb}

\title{Gated Recurrent Memory for Long-Horizon Time-Series Forecasting}
\author{R. Petrova \and J. Owusu \and L. Beaumont}
\date{}

\begin{document}
\maketitle

\begin{abstract}
We study LSTM-based forecasters for hourly load and meteorological signals
spanning multi-year horizons. A two-layer stacked LSTM with peephole
connections, trained on sliding windows with a curriculum on horizon length,
reduces forecast MAPE by $18\%$ over a strong seasonal-naive baseline at the
seven-day horizon. The key trick is not architectural but procedural: we
match teacher forcing rates to the empirical signal-to-noise ratio of each
horizon.
\end{abstract}

\section{Problem}
Given a multivariate series $x_{1:T}$, predict $\hat{x}_{T+1:T+H}$ for
horizons $H \in \{24, 168, 720\}$ hours. The two challenges are accumulating
errors past about three days and capturing exogenous regimes (holidays,
weather fronts) that the input window does not see directly.

\section{Model}
The recurrent cell maintains hidden state $h_t$ and cell state $c_t$:
\begin{align}
i_t &= \sigma(W_i x_t + U_i h_{t-1} + V_i c_{t-1} + b_i), \\
f_t &= \sigma(W_f x_t + U_f h_{t-1} + V_f c_{t-1} + b_f), \\
o_t &= \sigma(W_o x_t + U_o h_{t-1} + V_o c_t + b_o), \\
c_t &= f_t \odot c_{t-1} + i_t \odot \tanh(W_c x_t + U_c h_{t-1} + b_c), \\
h_t &= o_t \odot \tanh(c_t).
\end{align}
Peephole terms $V_\bullet c_\bullet$ let the gates read the current cell state
directly, which we found materially helps preserve trend information across
the daily reset.

\section{Training}
We use the Adam optimiser with $\eta = 10^{-3}$ decayed by $0.5$ every twenty
epochs. The teacher-forcing ratio anneals from $0.9$ to $0.0$ over the first
sixty epochs. Crucially, we sample horizon length $h \sim \mathcal{U}(1, H)$
during training so the network never specialises to a single decoding depth.

\section{Datasets}
We evaluate on three corpora: (i) ENTSO-E hourly electricity load for ten
European countries, (ii) NOAA station-level surface temperature, (iii) the
MIMIC-III ICU vitals subset filtered to thirty-minute resampling. The
electricity series has the strongest seasonal pattern; the ICU vitals are
the most non-stationary.

\section{Results}
\begin{tabular}{lccc}
\hline
Model & MAPE@24h & MAPE@168h & MAPE@720h \\
\hline
Seasonal-naive & 6.2\% & 9.1\% & 12.4\% \\
Vanilla LSTM   & 5.1\% & 8.7\% & 13.0\% \\
Stacked + peephole & 4.4\% & 7.5\% & 11.1\% \\
\hline
\end{tabular}

The vanilla LSTM beats the baseline at short horizons but actually does
worse at one month --- a textbook case of overconfident extrapolation.
The peephole variant degrades more gracefully because the cell state, not
the noisy hidden output, carries the long-range trend.

\section{Discussion}
Time-series forecasting is one of the few domains where Transformers do not
yet decisively beat recurrent baselines. Our reading is that the inductive
bias of LSTMs --- a learned, write-controlled accumulator --- matches the
underlying generative process more directly than self-attention does for
slowly-evolving signals. Curriculum-on-horizon is a cheap addition that
should generalise to any sequence model.

\bibliographystyle{plain}
\begin{thebibliography}{9}
\bibitem{hochreiter1997long} S. Hochreiter, J. Schmidhuber. \emph{Long Short-Term Memory.} Neural Computation, 1997.
\bibitem{gers2002learning} F. Gers et al. \emph{Learning Precise Timing with LSTM Recurrent Networks.} JMLR, 2002.
\end{thebibliography}

\end{document}
